\section{La estructura de producto del modelo como OS: de la puerta de entrada del chat a la puerta de entrada de las tareas}

Llamar al modelo una nueva generación de OS es lo que más fácilmente se malinterpreta como una exageración. Dicho con más precisión, el modelo podría convertirse en la «puerta de entrada de las tareas»: el usuario ya no abre primero una app, sino que primero expresa un objetivo, y el modelo elige las herramientas, invoca los servicios, organiza el contexto, ejecuta las acciones y devuelve el resultado. Este cambio reconfigurará la distribución de aplicaciones, la forma del SaaS y la lógica de compra del software empresarial.

\subsection{Estructura de producto en tres capas}

\noindent\begin{longtable}{p{0.22\textwidth}p{0.34\textwidth}p{0.35\textwidth}}
\toprule
Capa & Responsabilidad & Foco de la competencia \\
\midrule
Capa del modelo & Comprender la intención, planificar tareas, generar código o acciones & Capacidad base, capacidad de Coding, capacidad de uso de herramientas. \\
Capa de Harness & Conectar aplicaciones, permisos, memoria, registros, reintentos, evaluación & Confiabilidad del flujo de trabajo, costo, seguridad y observabilidad. \\
Capa de aplicación & Alojar los objetos de negocio concretos y la relación con el usuario & Datos, escenarios, distribución, hábitos del usuario y pago. \\
\bottomrule
\end{longtable}

\begin{knowledgebox}{Por qué la conversión en OS amenaza la puerta de entrada de las aplicaciones}
Si el usuario expresa sus objetivos a través del modelo, y el modelo se encarga de elegir las aplicaciones e invocar los servicios, el valor de las apps tradicionales como puerta de entrada disminuirá. Las aplicaciones siguen siendo importantes, pero podrían pasar de «el usuario las abre por iniciativa propia» a «el modelo las asigna según la tarea».
\end{knowledgebox}

\subsection{Los cambios en el software empresarial}

El software empresarial se diseñaba en torno a interfaces operadas por personas: formularios, botones, flujos, permisos e informes. En la era de los Agent, el software empresarial también necesita ofrecer API que el modelo pueda invocar, un estado legible, registros con capacidad de auditoría y operaciones recuperables. Quien convierta antes su producto en un agent-friendly workflow tendrá más facilidad para entrar en la red de asignación de tareas del modelo como OS.

\begin{importantbox}{Características del software agent-friendly}
API claras, permisos explícitos, estado verificable, acciones que se pueden revertir, registros estructurados, sandbox de bajo costo y buena documentación. Sin esto, por muy capaz que sea el modelo, difícilmente podrá operar los sistemas empresariales de forma segura y estable.
\end{importantbox}

\subsection{Resumen de la sección}

El modelo como OS no sustituye al núcleo del sistema operativo, sino que se convierte en la capa de asignación de tareas. La competencia futura del software pasará de «quién controla la puerta de entrada de la interfaz» a «quién puede ser invocado de forma confiable por el modelo».
